\section*{September 8, 2026: Remove obsolete parcel records and locate individual homes}
\textit{Agent-drafted follow-up. Jacob authorized supported corrections and requested a separate discussion of questionable cases.}

A general rule now removes 36 obsolete single-home records while preserving
all 36 corresponding homes under their newer parcel numbers. It requires a
unique point match within one foot, matching recorded street addresses,
land areas within 0.5 percent, consecutive assessment histories, and a successor
construction year no later than the first old assessment. This last restriction
helps distinguish administrative renumbering from a later replacement building.
Twenty-five matches are the Morgan homes identified in the preceding review;
eleven are elsewhere. None of the 36 old identifiers appears in the frozen
paper input. Their retained successors' years, units, floor areas and lot areas
are unchanged.

Maplewood's old ten-card record is also redundant. The newly retrieved exact
2018 parcel contains precisely the ten retained individual homes. Each home
reports 2,060 building square feet, matching every old card. The map measures
14,310.04 square feet, agreeing with the old assessment's 14,310. The existing
two-year, two-percent complete-home matching rule now retires that aggregate.
All ten individual measurements and years remain unchanged. This uses a pinned
historical polygon and the general rule, rather than a new manual exclusion.
The aggregate and ten individual identifiers occur in the frozen paper data,
all outside 500 feet; that snapshot has not been changed.

Green Street's mistyped parcel reference is corrected in source preparation,
with the original reference preserved. The resulting two-parcel building has
six units, 9,036 building square feet and 6,818 Assessor land square feet.
The 2013 polygons independently total 6,818.78 square feet. Its historical
parcel-center distance is approximately 783 feet. One obsolete source group
disappears. This is the only row-specific correction adopted in this work;
the committed ledger identifies the exact assessment row and erroneous reference.

The location producer now distinguishes individual homes from their former
large development parcels. The approved next-year exact-parcel coordinates
supply 1,341 individual locations inside the historical polygons. Relative to
the previous reconstructed data, 13 move inside 500 feet and 37 move outside.
The five King/Calumet/Eastgate distances become approximately 597, 616, 615,
420 and 313 feet. Individual Assessor land measurements are preserved; the
former site's polygon is coverage evidence, not each home's denominator.
Another 2,113 former-parcel centers lack an accepted individual point and are
explicitly labeled accordingly. Complete parcel coverage does not resolve those
individual-location questions. No coordinates beyond the approved one-year
window were adopted.

Lake Park remains a proposed decision. Four 2016 permits identify four five-home
townhouse buildings. Twenty current individual records report construction in
2019, while twenty old individual records report 2014 or 2017 and hand off to
the current records after 2019. A ten-card parent also needs its membership
reconciled. The recommendation is to preserve twenty individual homes and accept
the later Assessor proxy, rather than retain old records as extra construction.
No Lake Park exclusion or year decision was applied. Deming still lacks
individual lot denominators, and the 64th Street review now requires reconciling
four nearby identifiers rather than assuming the initially suggested pair.

\clearpage
\textit{Reproduction and limits:} Source preparation, pinned parcel acquisition,
Assessor project construction and historical boundary assignment build through
their task Makefiles. Standard reports describe the saved products; the audit
independently recomputes distances to every boundary of the assigned ward.
An unchanged second build runs no substantive producer. Removing the generated
point file and rebuilding with two Make jobs reproduces all point attributes
and geometry exactly. The Maplewood addition preserves all 12,987 previously pinned parcel features'
attributes and geometry. The corrected history retains every original parcel
reference and changes exactly one. The working branch is
\texttt{sales\_sample\_exploration}, following revision \texttt{f8b28408}.
Earlier unrelated logbook products are held at their existing versions solely
for document compilation. The paper's 8,648-row input remains unchanged;
no density regression was rerun. This is progress on the reconstruction, not
a claim that final density eligibility or full replication is complete.

The saved pre-change preferred-candidate file provides the controlled comparison:
only Green Street's housing and land measurements change. Thirty-six old-home
records and the Maplewood aggregate change to duplicate exclusions; Green Street
changes from review to retention. One Green Street source identifier disappears.
Several older intermediate reports were last recorded before subsequent
reconstruction work. Their refreshed counts, including 9,465 exact permit links
versus 9,356 in the old report, are not a controlled estimate of this batch's
changes. The current producers and their key checks passed, but no pre-change
copy of each such intermediate was saved for a row-by-row comparison.
